% Generated by `taskgen course examples` from powmod.addition_chain (seed 3).
% Do not edit: change examples.yaml and regenerate.
\begin{example}[аддитивная цепочка]\label{ex:exp-chain}
Построить аддитивную цепочку для $e = 55$, соответствующую бинарному алгоритму, и найти более короткую аддитивную цепочку. С её помощью вычислить \[55^{55}\ (\textrm{mod}\ 377).\]

\par\noindent\textit{Решение.} $e = 110111_2$. Бинарный алгоритм (слева направо: удвоение, при бите $1$ ещё прибавление единицы) даёт цепочку \[1,\ 2,\ 3,\ 6,\ 12,\ 13,\ 26,\ 27,\ 54,\ 55\]длины $\lambda(e) + \nu(e) - 1 = 5 + 5 - 1 = 9$ ($\lambda(e) = \lfloor \log_2 e \rfloor$, $\nu(e)$ --- число единиц). Более короткая цепочка длины $8$:\[1,\ 2,\ 4,\ 8,\ 16,\ 18,\ 36,\ 54,\ 55.\]Она кратчайшая: нижняя оценка $\ell(e) \ge \lceil \log_2 e \rceil = 6$, а полный перебор цепочек длины $7$ показывает, что их не существует. Вычисления по цепочке ($c_i = c_j + c_k$, $a^{c_i} = a^{c_j} a^{c_k}$):\[\begin{array}{c|c|r}i & c_i = c_j + c_k & a^{c_i} \bmod 377 \\ \hline 0 & 1 & 55 \\ 1 & 2 = 1 + 1\ (\text{кв.}) & 9 \\ 2 & 4 = 2 + 2\ (\text{кв.}) & 81 \\ 3 & 8 = 4 + 4\ (\text{кв.}) & 152 \\ 4 & 16 = 8 + 8\ (\text{кв.}) & 107 \\ 5 & 18 = 16 + 2 & 209 \\ 6 & 36 = 18 + 18\ (\text{кв.}) & 326 \\ 7 & 54 = 36 + 18 & 274 \\ 8 & 55 = 54 + 1 & 367 \\ \end{array}\]\[55^{55} \equiv 367\ (\textrm{mod}\ 377).\]
\end{example}
